\section{Missões interplanetárias}

Para expandir a presença humana além da Terra, outras missões estão em fase de pesquisa e desenvolvimento para possibilitar a colonização de outros planetas, como Marte.

\subsection{Mars Sample Return}

A missão \emph{Mars Sample Return (MSR)} \cite{nasa_mars_sample_return_2025} é uma iniciativa conjunta da NASA e da Agência Espacial Europeia (ESA) que deseja coletar e trazer para a Terra amostras de rochas e do solo marcianos para realizar análise de suas propriedades em laboratórios terrestres.
Com o intuito de preparar o caminho para futuros exploradores humanos em Marte e de aprofundar o entendimento sobre Marte e o sistema solar, essa campanha é considerada uma das mais desafiadoras já planejadas.

O conceito da missão envolve várias etapas.
Inicialmente, o rover \emph{Perseverance}, que já está em Marte, coleta e armazena amostras de rochas e solo, e os armazenam em tubos selados.
Depois disso, uma futura missão enviará o \emph{Sample Retrieval Lander}, que pousará próximo ao rover, na cratera Jezero, para receber as amostras coletadas. 
Este módulo de aterrissagem contará com um pequeno foguete, no qual as amostras serão carregadas para serem lançadas na órbita de Marte.
Uma vez em órbita, os \emph{Mars Ascent Vehicle} (e as amostras) serão capturadas pelo \emph{Earth Return Orbiter} \cite{nasa_earth_return_orbiter_2025}, uma espaçonave projetada para trazê-las de volta à Terra para análises aprofundadas.

A coleta e o retorno dessas amostras marcianas são fundamentais para entender sobre a possibilidade de vida em Marte, a história geológica do planeta e sua evolução climática.
As amostras retornadas permitirão realizar estudos e realizar descobertas que não seriam viáveis apenas com os dados coletados por sondas e pelo rover em Marte.
% Além disso, existe a exploração de asteroides, onde sistemas de navegação autônomos realizam aproximações com corpos irregulares e de baixa gravidade.

\subsection{Navegação autônoma na proximidade de asteroides}

A navegação autônoma de espaçonaves na proximidade de asteroides próximos à Terra (\emph{Near Earth Asteroid, NEO}) possui complexidade adicional envolvida por conta da estimativa de suas trajetórias.
Para tanto, se faz necessário a utilização de múltiplos sensores embarcados, como câmeras, sensores de atitude e LiDAR, ainda como explica \cite{vetrisano2015autonomous}, a fusão de dados de medições interespaciais com as estimativas calculadas de posição relativas ao asteroide pode reduzir a imprecisão nos dados.

O uso de medições relativas entre as espaçonaves, durante um voo de formação, pode contribuir para a mitigação de erros associados aos modelos de perturbações gravitacionais e variações no formato e rotação dos asteroides.
Ao combinar as medições de linha de visada com sensores baseados no efeito Doppler do Sol, conforme explica \cite{yim2000autonomous}, podem determinar a órbita do asteroide com maior precisão.
O voo em formações de espaçonaves que orbitam um asteroide requer controle autônomo devido à natureza não linear e imprevisível da dinâmica orbital dos asteroides, além disso, a radiação solar pode influenciar negativamente na estabilidade da órbita, e irá exigir técnicas mais avançadas de controle para manter a formação desejada \cite{vetrisano2015autonomous}.